\begin{proof}
\begin{enumerate}[label={$\langle1\rangle$\arabic*.},leftmargin=*]
\item 
\noindent\textbf{Step 1.} The equation of $\mathcal C$ is satisfied by both sections; the square root
is holomorphic for $|t|$ small, and the two sections are disjoint because
$O$ has $z=0$ whereas $P$ has $z=1$.  On the normalization of $C_0$, write
$r=u/v$.  Then $O(0)$ corresponds to $r=\infty$, $P(0)$ to $r=2$, and the
two branches over the node correspond to $r=1$ and $r=-1$.  Hence
\[
 \nu^*(B|_{C_0})\simeq\cO_{\PP^1}([2]-[\infty]).
\]
The rational function $f(r)=r-2$ has precisely this divisor and trivializes
the pullback.  Its values at the two branches are $f(1)=-1$ and
$f(-1)=-3$, so the induced gluing multiplier is $f(1)/f(-1)=1/3$,
up to inversion.  This is not a root of unity, and the lemma follows from
the gluing description in Lemma~\ref{wbII:lem:gluing-bundle}.

\end{enumerate}
\end{proof}

\begin{proof}
\begin{enumerate}[label={$\langle1\rangle$\arabic*.},leftmargin=*]
\item 
Because $S$ is the principal fibre $(t=0)$, its conormal line is trivial:
\(N^*_{S/\mathcal X}\simeq\cO_S\).  The projection
\(\operatorname{pr}_1:S\to C_0\) has
\(\operatorname{pr}_{1*}\cO_S\simeq\cO_{C_0}\), so \eqref{wbII:eq:CN:4.4} gives
\[
 H^0(S,A)\simeq H^0(C_0,A_\lambda)=0.
 \tag{W.165}\hypertarget{tag-03E2}{}\label{wbII:eq:CN:4.7}
\]
This proves the first vanishing in \eqref{wbII:eq:CN:4.5}.

We next check the second.  At a normal-crossing (self-crossing) point with local equation
\(uv=0\), the pullback map
\(\Omega^1_S\to\eta_*\Omega^1_{\widetilde S}\) has kernel exactly
\(\mathscr T_S\).  Indeed, in
\(R=\C\{u,v,w\}/(uv)\) the relation is
\(v\,du+u\,dv=0\).  If
\(a\,du+b\,dv+c\,dw\) restricts to zero on both branches, then
\(a\in(v)\), \(b\in(u)\), and
\(c\in(u)\cap(v)=(uv)=0\).  Modulo the relation the class is therefore a
multiple of $\tau=v\,du=-u\,dv$.  This class is nonzero: if
\(v\,du=h(v\,du+u\,dv)\), then coefficient comparison gives
  \((1-h)v=0\) and $hu=0$.  Since
\(\operatorname{Ann}_R(v)=(u)\) and \(\operatorname{Ann}_R(u)=(v)\),
\(1-h\in(u)\) and \(h\in(v)\), which would imply
\(1=(1-h)+h\in(u)+(v)\), impossible because $(u)+(v)$ is proper.
Moreover $(u+v)\tau=0$.  Since $u+v$ is a non-zero-divisor in
$R$ (the associated primes are $(u)$ and $(v)$), $\tau$ is genuine
torsion.
Conversely, each summand of the normalization module is a module over one of
the domains $R/(u)$ and $R/(v)$.  A non-zero-divisor of the reduced ring $R$
avoids both minimal primes and therefore acts injectively on both summands;
the target is torsion-free.  Thus the kernel of the normalization map is
exactly the torsion subsheaf, and there is an injection
\[
 \widetilde\Omega^1_S\hookrightarrow\eta_*\Omega^1_{\PP^1\times\PP^1}.
 \tag{W.166}\hypertarget{tag-03E3}{}\label{wbII:eq:CN:4.8}
\]
Since $\eta^*A\simeq\cO_{\PP^1\times\PP^1}$, the projection formula and \eqref{wbII:eq:CN:4.8}
give
\[
 H^0(S,\widetilde\Omega^1_S\otimes A)
 \hookrightarrow
 H^0(\PP^1\times\PP^1,\Omega^1_{\PP^1\times\PP^1})=0.
 \tag{W.167}\hypertarget{tag-03E4}{}\label{wbII:eq:CN:4.9}
\]
The last equality follows from
\(\Omega^1_{\PP^1\times\PP^1}\simeq
\operatorname{pr}_1^*\Omega^1_{\PP^1}\oplus
\operatorname{pr}_2^*\Omega^1_{\PP^1}\) and
\(H^0(\PP^1,\Omega^1_{\PP^1})=0\).

Finally, along the conductor $D$ the pullback of $dg|_S$ is zero.  Locally
$g=h\,uv$ with $h$ a unit, so after restriction to $S$ one has
$i^*(dg)=h(v\,du+u\,dv)$; both coefficients vanish on the preimage of $D$.
Proposition~\ref{wbII:prop:conductor}
therefore gives a nonzero section in \eqref{wbII:eq:CN:4.6}.  This proves the theorem.

\end{enumerate}
\end{proof}
